\section{Discussion}
\label{sec:discussion}

\subsection{Why reshaping spectra into images improves trait retrieval}
\label{sec:disc_reshape}

% --- Opening: headline finding in context ---
 - This study demonstrates that a simple reshape of 1D
   hyperspectral reflectance into a 2D image significantly
   improves multi-trait retrieval ($R^2 = 0.684$ vs 0.587
   for supervised 1D-CNN)
 - Consistent across all eight traits --- no trait is
   degraded by the 2D representation
 - Reshape outperformed five alternative transformations,
   confirming prior findings that direct reshaping is
   competitive or superior to engineered encodings
   \citep{hennessy2022reshaping, deev2024spectrum,
   sharma2019deepinsight}, and extending this to
   multi-trait regression with heterogeneous field data

% --- Mechanism: multi-scale receptive field ---
 - The $42 \times 42$ reshape creates a topology where a
   single $3 \times 3$ kernel spans $\sim$84\,nm of
   spectral range (vertical stride of 42\,nm), versus
   $\sim$3\,nm for an equivalent 1D kernel
 - After a few pooling stages in EfficientNet-B0, the
   effective receptive field covers hundreds of nm ---
   enough to simultaneously capture pigment absorption
   at 680\,nm, the red-edge at 750\,nm, and SWIR
   structural features
 - In a 1D-CNN, achieving the same spectral reach
   requires very deep networks or large kernels ($k > 80$),
   increasing overfitting risk with limited labeled data
   ($n = 4{,}508$) \citep{cherif2023spectra}
 - The reshaped image also encodes trait-specific
   ``spectral textures'' --- visually distinct spatial
   patterns arising from different biochemical
   compositions (Figure~\ref{fig:transforms_example}) ---
   that 2D convolutions are specifically designed to
   detect

% --- Per-trait evidence supports the mechanism ---
 - The per-trait improvement gradient directly supports
   the cross-band correlation mechanism
 - Largest gains for traits with broad, distributed
   absorption features:
   C$_\text{bc}$ ($\Delta R^2 = +0.245$, carbon-based
   constituents absorbing across 1700--2300\,nm),
   Anth ($\Delta = +0.174$, interacting with chlorophyll
   features $\sim$80--180\,nm apart),
   and C$_\text{m}$ ($\Delta = +0.122$, dry matter
   spanning multiple SWIR windows)
   \citep{curran1989imaging, feret2021prospectpro}
 - Moderate gains for pigments (C$_\text{ab}$
   $\Delta = +0.080$, C$_\text{ar}$ $\Delta = +0.070$):
   narrower visible-range features already partially
   captured by 1D convolutions, though 2D adds
   correlations with red-edge and NIR regions
   \citep{homolova2013review}
 - Near-zero gain for C$_\text{p}$ ($\Delta = +0.008$):
   protein has a narrow absorption at $\sim$2100\,nm
   \citep{feret2021prospectpro} --- cross-band
   correlations add minimal information for spectrally
   confined signals
 - This gradient (broad features $\rightarrow$ largest
   gains; narrow features $\rightarrow$ smallest gains)
   is consistent with the hypothesis that the 2D
   advantage arises from multi-scale cross-band
   integration

% --- Why complex transforms underperform ---
 - CWT ($R^2 = 0.641$) and Spectrogram ($R^2 = 0.636$)
   formed a second tier --- both preserve some spectral
   locality through windowed decomposition
   \citep{addison2017wavelet, shuai2025multichannel},
   but placing wavelength on one axis and scale/frequency
   on the other breaks the wavelength-to-wavelength
   adjacency that Reshape preserves
 - These transforms have proven effective for
   classification tasks where scale-specific
   discriminative features matter
   \citep{ong2025sugarcane, zhuang2023cwt,
   shuai2025multichannel, zhang2025gaf}, but for
   continuous regression the direct wavelength topology
   appears more informative
 - GAF ($R^2 = 0.609$) maps spectral values to angular
   space, grouping bands by reflectance magnitude rather
   than wavelength position, thus losing the physical
   ordering that connects absorption features to
   biochemistry \citep{wang2015gaf, ganesan2021gaf};
   GAF and MTF were designed for time-series
   classification where phase dynamics carry
   discriminative information --- a fundamentally
   different task
 - 2D-COS ($R^2 = 0.555$, std $= 0.077$) showed the
   highest seed variability; our single-sample adaptation
   via windowed segmentation \citep{noda2004twodimensional}
   is an ad-hoc perturbation that may introduce noise
   rather than meaningful correlation structure
 - MTF ($R^2 = 0.419$) was the only transform worse than
   the 1D baseline --- quantile-based discretisation
   destroys continuous spectral information critical for
   regression \citep{wang2015gaf}
 - Overall: transforms preserving natural wavelength
   ordering and continuous values outperform those
   remapping spectral information into abstract
   mathematical spaces


\subsection{\jl{What the model learns from the 2D representation}}
\label{sec:disc_importance}

% --- 2D importance makes the receptive-field mechanism empirical ---
\jl{ - The variable importance analysis gives empirical
   support to the receptive-field mechanism proposed in
   Section~\ref{sec:disc_reshape}. Grad-CAM relevance
   forms compact 2D clusters. Once unfolded to wavelength,
   six of the eight traits concentrate importance in known
   absorption regions above chance
   (Figure~\ref{fig:importance}).}

\jl{ - The unfolded maps expose how the model reads the
   image. Horizontal neighbours are adjacent bands.
   Vertical neighbours are bands 42 positions apart. A
   single 2D kernel reads both at once. The model can
   co-weight a visible feature and its shortwave-infrared
   overtone about 42 bands away inside one receptive
   field. A 1D kernel of the same size cannot.}

\jl{ - Integrated Gradients and Grad-CAM agree. IG works at
   the fine, signed, band level. Grad-CAM works at the
   coarse, region level. Their agreement raises confidence
   that the importance reflects learned spectral structure
   rather than an artifact of one method.}

% --- Relation to known absorption regions and why they differ ---
\jl{ - Six traits concentrate importance in physically
   grounded regions. Anthocyanins, leaf mass per area,
   carbon-based constituents, chlorophyll, and protein all
   stay above chance. Their peaks fall on the red-edge
   (C\textsubscript{ab} at 703\,nm) and on shortwave-infrared
   cellulose, lignin, and protein features (near 2144\,nm).
   These match constituent absorption reported at the leaf
   level \citep{serbin2020scaling, angel2022remote,
   curran1989imaging, feret2021prospectpro}.}

\jl{ - These windows come from constituent absorption
   physics, not from the model. Pigment electronic
   transitions in the visible, and water, protein,
   cellulose, and lignin vibrational overtones in the
   shortwave infrared, occupy stable wavelength ranges
   across studies \citep{curran1989foliar, fourty1996leaf,
   feret2021prospectpro}. They are independent of the
   retrieval method. Band importance derived from the model
   itself shifts with the algorithm, the regularisation,
   and the training data. We compared against the physical
   regions to keep the reference independent of our model.}

\jl{ - Carotenoids ($0.82\times$) and equivalent water
   thickness ($0.81\times$) fall below chance. The model
   predicts both well, so the mismatch reflects where the
   model looks rather than how well it performs. Several
   non-exclusive reasons explain the difference from the
   literature.}

\jl{ - Scale differs. The reference regions come from leaf
   or constituent spectra. GreenHyperSpectra holds canopy
   reflectance. Canopy structure, leaf area index, soil
   background, and multiple scattering broaden, shift, and
   saturate absorption features. The predictive bands need
   not match the leaf-level maxima \citep{homolova2013review}.}

\jl{ - Traits covary. The multi-trait model can predict a
   weakly featured trait through a correlated strong one.
   Carotenoids share variance with chlorophyll. The model
   leans on the chlorophyll red-edge at 703\,nm rather than
   the weak carotenoid feature at 440--520\,nm. A good
   R\textsuperscript{2} of 0.707 with low enrichment reflects
   this shortcut.}

\jl{ - Preprocessing removed water bands. The strongest
   water features near 1450 and 1940\,nm fall partly inside
   the removed gaps (1351--1430 and 1801--2050\,nm). The
   model uses the surviving edge at 1431\,nm. The low water
   enrichment is partly a preprocessing artifact, not
   biology.}

\jl{ - The metric depends on our window choice. We set the
   absorption windows from the literature. Narrow windows
   penalise the legitimate use of adjacent bands and lower
   the enrichment for traits with broad or shifted
   features.}


\subsection{Self-supervised pretraining in the 2D framework}
\label{sec:disc_pretraining}

% --- Key comparison ---
 - The 2D supervised baseline ($R^2 = 0.684$) surpassed
   all 1D methods, including the best self-supervised
   MAE-1D FT ($R^2 = 0.645$) of
   \citet{cherif2025greenhyperspectra}
 - The representational advantage of 2D images alone
   (+0.039 over MAE-1D FT) exceeds the benefit of
   pretraining on 139K unlabeled spectra within the 1D
   framework (+0.058 over supervised 1D)
 - The 2D representation is therefore the dominant factor,
   not the learning strategy

% --- MAE-2D results ---
 - MAE-2D FT ($R^2 = 0.667 \pm 0.006$) outperformed
   MAE-1D FT by +0.022, confirming the 2D advantage
   extends to self-supervised learning
 - However, MAE-2D FT did not surpass supervised 2D
   ($-0.017$), contrasting with natural image domains
   where MAE pretraining typically helps \citep{he2022mae}
 - Linear probing ($R^2 = 0.646 \pm 0.020$, only 38.6K
   trainable parameters) exceeded all 1D baselines
   including MAE-1D FT --- capturing 97\% of fine-tuned
   performance and indicating the frozen encoder already
   encodes most trait-relevant spectral information

% --- Why supervised 2D > MAE-2D FT ---
 - Three non-mutually exclusive hypotheses:
   (1) Architecture mismatch --- EfficientNet-B0's spatial
   inductive biases (local connectivity, translation
   equivariance) outperform ViT's learned attention with
   only $n = 4{,}508$ labeled samples
   \citep{dosovitskiy2021vit};
   (2) Easy reconstruction --- reshaped spectral images
   have smooth, continuous structure unlike natural images,
   so the MAE task converges rapidly ($\sim$10 epochs,
   val\_loss $\sim$0.0001) without learning deep
   trait-discriminative features;
   (3) Sufficient labeled data --- 4,508 samples spanning
   diverse ecosystems may already capture most learnable
   spectral--trait relationships
 - Practical implication: researchers can achieve
   state-of-the-art performance with reshape + standard
   CNN, without the computational cost of pretraining on
   large unlabeled corpora
 - The MAE-2D result remains valuable for data-scarce
   settings where few labeled samples are available


\subsection{Implications for trait-based ecology and remote sensing}
\label{sec:disc_implications}

% --- Carbon cycle and ecosystem functioning ---
 - The largest improvement was for C$_\text{bc}$
   ($\Delta R^2 = +0.245$), directly linked to carbon
   stocks, litter decomposition, and nutrient cycling
   \citep{kattge2020try}
 - Improved C$_\text{bc}$ and C$_\text{m}$ retrieval
   supports non-destructive estimation of leaf dry
   matter investment --- a key axis of the leaf economics
   spectrum \citep{wright2004les} --- enabling more
   reliable mapping of ecosystem carbon pools and
   decomposition dynamics for Earth system models
   \citep{homolova2013review}

% --- Functional diversity monitoring ---
 - Multi-trait retrieval from a single model preserves
   inter-trait covariance structure, essential for
   functional diversity metrics (FDis, FRic, FEve, Rao's
   Q) \citep{villeger2008fd, laliberte2010fdis}
 - Trait-by-trait retrieval with independent models may
   distort correlations and bias diversity estimates
 - Simultaneous improvement across all eight traits means
   functional trait spaces from 2D predictions will be
   more accurate and internally consistent than 1D
   baselines --- relevant for biodiversity monitoring
   programmes relying on spectral trait mapping
   \citep{schweiger2018plant, binh2025mangrove, weiss2020remote}

% --- Operational simplicity ---
 - Reshape requires no domain-specific tuning (no wavelet
   scales, STFT windows, quantile bins, or perturbation
   parameters) and is implementable in $\sim$6 lines of
   code
 - Robustness across seeds (std $= 0.001$) ensures
   reproducible results without hyperparameter sensitivity
 - This lowers the barrier for ecologists and remote
   sensing practitioners who need state-of-the-art
   predictions without signal processing expertise

% --- Scaling to hyperspectral missions ---
 - Upcoming and recently launched missions --- EnMAP
   \citep{guanter2015enmap}, PRISMA
   \citep{verrelst2021prisma}, SBG/EMIT
   \citep{cawse2021sbg, green2020emit} --- will produce continental-
   to global-scale hyperspectral archives with contiguous
   spectra across 400--2500\,nm
 - A simple, robust method validated across 50
   heterogeneous datasets is well suited for operational
   deployment on these data streams

% --- Data-efficient learning for undersampled biomes ---
 - The MAE-2D linear probe ($R^2 = 0.646$, 38.6K params)
   opens a pathway for trait prediction in data-scarce
   biomes (tropical forests, tundra, peatlands) that are
   underrepresented in existing trait databases
   \citep{kattge2020try, zhang2025ppadanet}
 - A frozen pretrained encoder requiring only a
   lightweight linear head could enable trait mapping
   with minimal local calibration data, aligning with the
   broader trend toward foundation models for remote
   sensing \citep{binfaruk2025terramae}


\subsection{Limitations and future directions}
\label{sec:disc_limitations}

% --- Methodological limitations ---
 - All transform comparisons used a single backbone
   (EfficientNet-B0); performance rankings could differ
   with other architectures, though
   \citet{cherif2023spectra} showed it to be optimal for
   this data scale
 - Fixed image size ($224 \times 224$): the reshape
   topology depends on resolution, and other sizes may
   favour different transforms --- a systematic study is
   warranted
 - Single-sample 2D-COS via windowed segmentation is an
   ad-hoc approximation of \citet{noda2004twodimensional}'s
   multi-sample formulation --- its poor performance may
   partly reflect this limitation
 - All experiments used GreenHyperSpectra; validation on
   independent datasets with different spectral
   resolutions, wavelength ranges, and ecosystem types is
   needed
 - The moderate labeled set ($n = 4{,}508$) may limit the
   benefit of self-supervised pretraining; with larger
   datasets, the gap between supervised and MAE-2D could
   narrow
 - This study used point-level spectra without spatial
   context --- airborne/satellite imagery introduces mixed
   pixels and atmospheric effects not addressed here

\jl{ - The Reshape geometry is arbitrary, which limits the
   importance maps. The model zero-pads the spectrum and
   resizes it bilinearly from $42 \times 42$ to
   $224 \times 224$, smearing each band over about five
   pixels. The vertical 42-band adjacency comes from the
   $\text{side} = 42$ fold, not from physics, so part of
   the cross-row structure reflects the layout. The
   unfolded importance has limited spectral resolution and
   the apparent peaks can shift by a few bands. Row-wrap
   boundaries place spectrally adjacent bands far apart in
   the image. We therefore read the importance maps as
   approximate.}

% --- Future directions ---
 - Investigate the effect of image size and reshape
   topology on retrieval --- different row widths create
   different cross-band neighborhoods that may favour
   different traits
 - Explore harder MAE pretraining: higher mask ratios,
   structured masking of spectral regions, and CWT
   scalograms as inputs (richer spatial structure than
   reshape)
 - Pretrain larger 2D encoders on massive spectral
   archives for general-purpose spectral feature
   extraction (foundation models)
 - Combine reshaped spectral images with spatial context
   from imaging spectrometers or LiDAR-derived canopy
   structure (multi-modal fusion)

